\section{おわりに}
本稿では，検出と命名を分離した粒度可変物体アノテーションに対し，SigLIPとシグモイドによる確信度，YOLOとVLMのスコアの比較による深さ0の決定，およびWordNetにおいてCOCOのカテゴリと同じ深さにある語彙を深さ0の候補とする3つの変更を，現在の構成と比較した．
シグモイドは，深さ0が誤った枝の識別性能をハイブリッド構成で0.418から0.617へ改善したが，深さ1の正誤の識別ではソフトマックスに劣った．
スコアの比較による切替は，YOLOとVLMの尺度の違いにより切替の頻度が大きく変わり，ソフトマックスではCOCOの精度を低下させた．
WordNetの同じ深さの語彙は，COCO外の物体の深さ0の正解率を3.8\%に上げる一方，COCOのTop-1正解率を0.945へ低下させた．
今後の課題として，ソフトマックスとシグモイドを組み合わせた確信度，YOLOとVLMの確信度の較正，および物体の名前としての用いられやすさを考慮した深さ0の語彙の選択が挙げられる．
